\documentclass[border=0pt]{standalone}
% The shared preamble of every docs diagram (docs/diagrams/<name>.tex).
% dvisvgm's driver first, so patterns and opacity survive into the SVG.
\def\pgfsysdriver{pgfsys-dvisvgm.def}
\usepackage[siunitx,RPvoltages]{circuitikz}
\usepackage{xcolor}
% Latin Modern: outline fonts for every glyph a label uses, µ included.
\usepackage[T1]{fontenc}
\usepackage{lmodern}
\usetikzlibrary{arrows.meta,backgrounds,calc,patterns,positioning}
\tikzset{
  % A white card, so diagrams read in mdBook's dark themes too.
  every picture/.append style={show background rectangle,
    background rectangle/.style={fill=white,inner sep=4pt}},
  % Layout geometry: layers told apart by fill and hatching, never colour alone.
  layer a/.style={draw=black,line width=0.6pt,fill=gray!20},
  layer b/.style={draw=black,line width=0.6pt,fill=gray!45},
  layer c/.style={draw=black,line width=0.6pt,fill=white,postaction={pattern=north east lines}},
  layer d/.style={draw=black,line width=0.6pt,fill=white,postaction={pattern=crosshatch dots}},
  marker/.style={draw=black,dashed,line width=0.6pt},
  % A measured distance and its label.
  dim/.style={{Latex[length=2mm]}-{Latex[length=2mm]},line width=0.5pt},
  dim label/.style={font=\small,fill=white,inner sep=1pt},
  % Where a rule fails: thick, dashed and labelled, not only red.
  violation/.style={draw=red!75!black,line width=1.4pt,dash pattern=on 3pt off 2pt},
  violation label/.style={font=\small\bfseries,text=red!75!black},
  caption/.style={font=\small},
}

\begin{document}
\begin{circuitikz}
  % Cross-section through the crossing: a and b on met1, c on met2 above, the substrate below.
  \fill[layer a] (-0.5,0) rectangle (7.5,-0.25);
  \draw[line width=0.6pt] (-0.5,0) -- (7.5,0);
  \node[caption,anchor=west] at (7.6,-0.12) {substrate (ground)};
  \fill[layer b] (1,2) rectangle (2,3);
  \fill[layer b] (4,2) rectangle (5,3);
  \fill[layer c] (0,3.8) rectangle (6,4.8);
  \node[caption,fill=white,inner sep=1pt] at (1.5,2.5) {a};
  \node[caption,fill=white,inner sep=1pt] at (4.5,2.5) {b};
  \node[caption,fill=white,inner sep=1pt] at (3,4.3) {c};
  \node[caption,anchor=west] at (6.1,2.5) {met1, $t = \SI{0.36}{\micro\metre}$};
  \node[caption,anchor=west] at (6.1,4.3) {met2};
  % The three closed forms.
  \draw (2,2.5) to[C] (4,2.5);
  \node[caption] at (3,3.2) {$\varepsilon_0\varepsilon_r t L / S$};
  \draw (1.5,3) to[C] (1.5,3.8);
  \node[caption,anchor=east] at (1.25,3.4) {$\varepsilon_0\varepsilon_r A / d$};
  \draw (4.5,3) to[C] (4.5,3.8);
  \draw (1.5,2) to[C] (1.5,0.6) node[ground]{};
  \node[caption,anchor=east] at (1.25,1.3) {$C_a A + C_f P$};
  \draw (4.5,2) to[C] (4.5,0.6) node[ground]{};
  \draw[dim] (2,1.6) -- (4,1.6) node[dim label,midway] {$S$};
  \draw[dim] (5.4,3) -- (5.4,3.8) node[dim label,midway,right=2pt] {$d$};
\end{circuitikz}
\end{document}
